%% Chapter 2 — Related Work
\label{ch:work}

This chapter reviews the work most relevant to the EMS Portal. It begins with
online service marketplaces and how they currently match supply with demand.
It then examines multi-criteria decision making, the formal basis for the
matching engine, and role-based access control, which underpins the platform's
authorization. The chapter concludes by identifying the gap this project
addresses.

\section{Current Advancements}

\subsection{Online Service Marketplaces}
Platforms for hiring technical providers fall into three groups. The first is
\textbf{horizontal generalist marketplaces} such as Upwork, Fiverr, and, in
Vietnam, vLance. They serve a broad range of work and rely on keyword search,
open bidding, and star ratings~\cite{kassi2019online}. A client describes a need
in plain language, the platform returns a list ranked mostly on text relevance
and popularity, and the client then weighs price, availability, location, and
reputation by hand. The second is \textbf{premium vetted networks} such as
Toptal and Arc.dev. They screen providers and then match them behind the scenes,
so the matching logic is not exposed to the client, and they focus on software
roles rather than broader engineering work. The third is \textbf{specialized
engineering platforms} such as Engre, Tasker, and Cad Crowd. Engre is a
directory of companies matched by static filtering and manual negotiation,
Tasker shortlists hardware specialists automatically but through a fixed,
black-box algorithm, and Cad Crowd runs design contests.

Table~\ref{tab:platforms} compares these platforms on the two capabilities
central to this work: an automatic multi-criteria ranking, and client- or
operator-\textit{configurable} weights. The comparison is based on each
platform's publicly available descriptions.

\begin{table}[ht]
	\centering
	\caption{How existing platforms match supply with demand.}
	\label{tab:platforms}
	\resizebox{\textwidth}{!}{%
		\begin{tabular}{llcc}
			\hline
			\textbf{Platform} & \textbf{Discovery / matching model} & \textbf{Auto multi-criteria} & \textbf{Configurable weights} \\ \hline
			Upwork / Fiverr & Keyword search, bidding, star ratings & \xmark & \xmark \\
			vLance (Vietnam) & Open bidding, manual short-listing & \xmark & \xmark \\
			Toptal & Vetting, then opaque behind-the-scenes match & \xmark & \xmark \\
			Arc.dev & AI-assisted match, software roles only & Partial & \xmark \\
			Engre & Company directory, static filter + manual & \xmark & \xmark \\
			Tasker & Automatic shortlist, fixed black-box algorithm & \cmark & \xmark \\
			\textbf{EMS Portal} & Weighted multi-criteria ranking, audited & \cmark & \cmark \\
			\hline
		\end{tabular}%
	}
\end{table}

While this model works well for a broad range of tasks, it falls short for
specialized engineering work. None of the surveyed platforms exposes
configurable matching weights, and only Tasker offers automatic multi-criteria
matching at all, but as a fixed black box the client can neither inspect nor
adjust the several competing factors that matter when selecting an
\textit{engineering} provider.

\subsection{Multi-Criteria Decision Making}
Picking a provider by several weighted factors is a case of
\textit{multi-criteria decision making} (MCDM). MCDM is a well-established field
with structured methods for ranking options that trade off competing
goals~\cite{triantaphyllou2000multi}. The simplest method is the
\textbf{weighted sum model}. Each option gets a score equal to the weighted sum
of its normalized factor values. More complex methods include the Analytic
Hierarchy Process, which sets factor weights from pairwise
comparisons~\cite{saaty1980analytic}. The weighted sum model fits an interactive
marketplace well. It is easy to understand. It is cheap to compute over many
candidates. And its weights can be tuned or compared without changing the
scoring logic. The matching engine relies on these properties.

\subsection{Role-Based Access Control}
Role-based access control (RBAC) gives permissions to roles, not to single
users. An authorization decision then follows from a user's role and the
permissions that role holds~\cite{sandhu1996role}. RBAC fits a marketplace with
a small, fixed set of actor types: client, service provider, and administrator.
It also fits a permission set defined per resource and action. The EMS Portal
uses this model directly. It encodes permissions in a canonical
\texttt{resource:verb:scope} form and maps roles to permission sets
(Section~\ref{sec:rbac-impl}).

\section{Research Gap}
Two gaps motivate this work.

\textbf{Gap 1: No Multi-Criteria Optimization in Provider Selection.}
As Table~\ref{tab:platforms} shows, none of the surveyed platforms, from
generalists like Upwork and vLance to specialized engineering platforms like
Engre and Tasker, treats provider selection as a multi-criteria optimization
problem the client can control. They offer no \textit{configurable} weighted
model to rank engineering providers by skill fit, availability, cost, proximity,
and rating simultaneously. Operators cannot compare a \textit{basic}
(equal-weight) policy against a \textit{tuned} alternative, and where automatic
matching exists at all, as in Tasker, it is a fixed black box.

\textbf{Gap 2: Marketplace Systems Lack Engineering Rigor.} Work in this space
is rarely designed or evaluated as a modular, testable, secure client-server
system with measured communication efficiency. Network performance is often
assumed rather than analyzed; security is treated as an afterthought.

The EMS Portal addresses both gaps. It implements a multi-criteria matching
engine with configurable weights stored as named policy objects, enabling
comparison of basic and tuned strategies. It also adopts a strict modular
architecture designed for security and testability, with explicit measurement of
network performance and functional correctness alongside matching effectiveness.
